\documentclass[11pt,a4paper]{article}

\usepackage[T1]{fontenc}
\usepackage[utf8]{inputenc}
\usepackage[polish]{babel}
\usepackage{lmodern}
\usepackage[margin=2.5cm]{geometry}
\usepackage{amsmath,amssymb}
\usepackage{booktabs}
\usepackage{array}
\usepackage{tabularx}
\usepackage{tikz}
\usepackage{caption}
\usepackage{enumitem}
\usepackage{xcolor}
\usepackage[hidelinks]{hyperref}

\newcommand{\todo}[1]{\textcolor{red}{\textbf{[TODO: #1]}}}

\title{Dron zwiadowczy w środowisku gridowym\\
  \large Projekt z przedmiotu \emph{Uczenie się ze wzmocnieniem}, 2026/2027}
\author{{Nicolas Liszek 179552}\\Politechnika Rzeszowska, WEiI}
\date{\today}

\begin{document}
\maketitle

\begin{abstract}
Projekt dotyczy agenta sterującego dronem zwiadowczym w dyskretnym środowisku
gridowym. Agent ma sfotografować zestaw celów i wrócić do bazy przy
ograniczonym zapasie energii, omijając strefy radarowe i ruchomy patrol.
Wybór kolejności celów wpływa na dalszy przebieg misji, ponieważ sfotografowanie
jednego z nich zamyka korytarz na mapie. Porównane zostaną agenty bazowe
(losowy i zachłanny) z agentami tablicowymi (Q-learning, SARSA, Q($\lambda$)).
\end{abstract}

\tableofcontents

\section{Opis problemu}
\label{sec:problem}

\paragraph{Repozytorium projektu.} Projekt jest realizowany w publicznym
repozytorium GitHub:
\url{https://github.com/nickx64/rl-recon-drone}.

Dron zwiadowczy startuje z bazy umieszczonej na dwuwymiarowej mapie
podzielonej na pola. Jego zadaniem jest wykonanie zdjęć $K$ celów
rozmieszczonych na mapie i powrót do bazy. Misję utrudniają:

\begin{itemize}
  \item \textbf{ograniczona bateria} -- każdy krok zużywa energię, a jej
        uzupełnienie jest możliwe wyłącznie na stacjach ładowania, które nie
        leżą na najkrótszych trasach między celami;
  \item \textbf{przeszkody statyczne} -- budynki, przez które nie można przelecieć;
  \item \textbf{strefy radarowe} -- przebywanie w nich skraca trasę, ale w
        każdym kroku grozi wykryciem drona i niepowodzeniem misji;
  \item \textbf{patrol} -- przeszkoda dynamiczna poruszająca się cyklicznie
        po stałej trasie; kontakt z patrolem kończy misję;
  \item \textbf{zdarzenie zależne od decyzji agenta} -- sfotografowanie
        celu nr~1 powoduje zamknięcie wybranego korytarza (pole $Z$), co może
        odciąć krótszą drogę do pozostałych celów.
\end{itemize}

Agent nie zna z góry modelu środowiska (prawdopodobieństwa wykrycia, skutków
zdarzeń) -- uczy się wyłącznie na podstawie interakcji. Projekt zakłada kilka
map o rosnącej trudności (\emph{easy}, \emph{medium}, \emph{hard}).
Przykładową mapę pokazano na rys.~\ref{fig:map-easy}.

\begin{figure}[h]
  \centering
  \begin{tikzpicture}[scale=0.55]
  \fill[blue!35] (0,9) rectangle ++(1,1);
  \node at (0.5,9.5) {\small\textbf{B}};
  \fill[white] (1,9) rectangle ++(1,1);
  \fill[white] (2,9) rectangle ++(1,1);
  \fill[white] (3,9) rectangle ++(1,1);
  \fill[black!70] (4,9) rectangle ++(1,1);
  \fill[white] (5,9) rectangle ++(1,1);
  \fill[white] (6,9) rectangle ++(1,1);
  \fill[white] (7,9) rectangle ++(1,1);
  \fill[white] (8,9) rectangle ++(1,1);
  \fill[yellow!50] (9,9) rectangle ++(1,1);
  \node at (9.5,9.5) {\small\textbf{2}};
  \fill[white] (0,8) rectangle ++(1,1);
  \fill[black!70] (1,8) rectangle ++(1,1);
  \fill[black!70] (2,8) rectangle ++(1,1);
  \fill[white] (3,8) rectangle ++(1,1);
  \fill[black!70] (4,8) rectangle ++(1,1);
  \fill[white] (5,8) rectangle ++(1,1);
  \fill[white] (6,8) rectangle ++(1,1);
  \fill[white] (7,8) rectangle ++(1,1);
  \fill[white] (8,8) rectangle ++(1,1);
  \fill[white] (9,8) rectangle ++(1,1);
  \fill[white] (0,7) rectangle ++(1,1);
  \fill[white] (1,7) rectangle ++(1,1);
  \fill[white] (2,7) rectangle ++(1,1);
  \fill[white] (3,7) rectangle ++(1,1);
  \fill[white] (4,7) rectangle ++(1,1);
  \fill[red!25] (5,7) rectangle ++(1,1);
  \fill[red!25] (6,7) rectangle ++(1,1);
  \fill[white] (7,7) rectangle ++(1,1);
  \fill[white] (8,7) rectangle ++(1,1);
  \fill[white] (9,7) rectangle ++(1,1);
  \fill[white] (0,6) rectangle ++(1,1);
  \fill[black!70] (1,6) rectangle ++(1,1);
  \fill[green!40] (2,6) rectangle ++(1,1);
  \node at (2.5,6.5) {\small\textbf{L}};
  \fill[white] (3,6) rectangle ++(1,1);
  \fill[black!70] (4,6) rectangle ++(1,1);
  \fill[black!70] (5,6) rectangle ++(1,1);
  \fill[black!70] (6,6) rectangle ++(1,1);
  \fill[orange!50] (7,6) rectangle ++(1,1);
  \node at (7.5,6.5) {\small\textbf{Z}};
  \fill[black!70] (8,6) rectangle ++(1,1);
  \fill[black!70] (9,6) rectangle ++(1,1);
  \fill[white] (0,5) rectangle ++(1,1);
  \fill[black!70] (1,5) rectangle ++(1,1);
  \fill[white] (2,5) rectangle ++(1,1);
  \fill[white] (3,5) rectangle ++(1,1);
  \fill[white] (4,5) rectangle ++(1,1);
  \fill[white] (5,5) rectangle ++(1,1);
  \fill[white] (6,5) rectangle ++(1,1);
  \fill[white] (7,5) rectangle ++(1,1);
  \fill[white] (8,5) rectangle ++(1,1);
  \fill[white] (9,5) rectangle ++(1,1);
  \fill[white] (0,4) rectangle ++(1,1);
  \fill[white] (1,4) rectangle ++(1,1);
  \fill[white] (2,4) rectangle ++(1,1);
  \fill[black!70] (3,4) rectangle ++(1,1);
  \fill[black!70] (4,4) rectangle ++(1,1);
  \fill[white] (5,4) rectangle ++(1,1);
  \fill[black!70] (6,4) rectangle ++(1,1);
  \fill[black!70] (7,4) rectangle ++(1,1);
  \fill[white] (8,4) rectangle ++(1,1);
  \fill[white] (9,4) rectangle ++(1,1);
  \fill[yellow!50] (0,3) rectangle ++(1,1);
  \node at (0.5,3.5) {\small\textbf{1}};
  \fill[white] (1,3) rectangle ++(1,1);
  \fill[white] (2,3) rectangle ++(1,1);
  \fill[white] (3,3) rectangle ++(1,1);
  \fill[black!70] (4,3) rectangle ++(1,1);
  \fill[white] (5,3) rectangle ++(1,1);
  \fill[white] (6,3) rectangle ++(1,1);
  \fill[black!70] (7,3) rectangle ++(1,1);
  \fill[white] (8,3) rectangle ++(1,1);
  \fill[green!40] (9,3) rectangle ++(1,1);
  \node at (9.5,3.5) {\small\textbf{L}};
  \fill[black!70] (0,2) rectangle ++(1,1);
  \fill[black!70] (1,2) rectangle ++(1,1);
  \fill[white] (2,2) rectangle ++(1,1);
  \fill[white] (3,2) rectangle ++(1,1);
  \fill[white] (4,2) rectangle ++(1,1);
  \fill[white] (5,2) rectangle ++(1,1);
  \fill[white] (6,2) rectangle ++(1,1);
  \fill[black!70] (7,2) rectangle ++(1,1);
  \fill[white] (8,2) rectangle ++(1,1);
  \fill[white] (9,2) rectangle ++(1,1);
  \fill[white] (0,1) rectangle ++(1,1);
  \fill[white] (1,1) rectangle ++(1,1);
  \fill[white] (2,1) rectangle ++(1,1);
  \fill[red!25] (3,1) rectangle ++(1,1);
  \fill[red!25] (4,1) rectangle ++(1,1);
  \fill[white] (5,1) rectangle ++(1,1);
  \fill[white] (6,1) rectangle ++(1,1);
  \fill[white] (7,1) rectangle ++(1,1);
  \fill[white] (8,1) rectangle ++(1,1);
  \fill[yellow!50] (9,1) rectangle ++(1,1);
  \node at (9.5,1.5) {\small\textbf{3}};
  \fill[white] (0,0) rectangle ++(1,1);
  \fill[white] (1,0) rectangle ++(1,1);
  \fill[white] (2,0) rectangle ++(1,1);
  \fill[red!25] (3,0) rectangle ++(1,1);
  \fill[red!25] (4,0) rectangle ++(1,1);
  \fill[white] (5,0) rectangle ++(1,1);
  \fill[black!70] (6,0) rectangle ++(1,1);
  \fill[white] (7,0) rectangle ++(1,1);
  \fill[white] (8,0) rectangle ++(1,1);
  \fill[white] (9,0) rectangle ++(1,1);
  \draw[step=1,gray!60] (0,0) grid (10,10);
  \draw[very thick,purple,dashed,<->] (5.5,5.5) -- (5.5,4.5) -- (5.5,3.5) -- (5.5,2.5);
\end{tikzpicture}

  \caption{Robocza mapa \emph{easy} ($10\times10$). B -- baza, 1--3 -- cele,
    L -- stacje ładowania, czerwone pola -- strefy radarowe, Z -- korytarz
    zamykany po sfotografowaniu celu~1, czarne pola -- przeszkody, fioletowa
    linia -- trasa patrolu (ruch tam i z powrotem).}
  \label{fig:map-easy}
\end{figure}

\paragraph{Dylemat zaprojektowany na mapie \emph{easy}.} Cel~1 leży najbliżej
bazy (6~kroków), ale jego sfotografowanie zamyka korytarz $Z$. Po tym
zdarzeniu jedyna droga do celu~2 prowadzi przez strefę radarową. Bezpieczna
kolejność to 2~$\rightarrow$~3~$\rightarrow$~1~$\rightarrow$~baza
(ok.\ 46 kroków po trasach omijających radary), co przekracza pojemność
baterii $E_{\max}=30$ -- agent musi więc dodatkowo zaplanować ładowanie.
Strategia zachłanna („najpierw najbliższy cel'') prowadzi natomiast do
konieczności lotu przez radar. Wartości zmierzono algorytmem BFS na mapie
z rys.~\ref{fig:map-easy}, bez uwzględnienia patrolu.

\paragraph{Cel agenta.} Wykonać zdjęcia wszystkich celów i wrócić do bazy
w możliwie małej liczbie kroków, nie dopuszczając do wykrycia ani do
rozładowania baterii.

\section{Uzasadnienie zastosowania uczenia się ze wzmocnieniem}
\label{sec:uzasadnienie}

Problem nie jest zadaniem klasyfikacji ani jednorazowej optymalizacji, ponieważ:

\begin{enumerate}
  \item \textbf{Decyzje są sekwencyjne} -- misja składa się z kilkudziesięciu
        kroków, a jakość pojedynczego ruchu zależy od całej dalszej trajektorii.
  \item \textbf{Decyzje zmieniają przyszły stan środowiska} -- zużycie energii
        ogranicza dalsze możliwości, a sfotografowanie celu~1 trwale zmienia
        mapę (zamknięcie korytarza $Z$). Kolejność odwiedzania celów ma więc
        znaczenie, a zachłanne wybieranie najbliższego celu może prowadzić
        do sytuacji bez wyjścia.
  \item \textbf{Występuje niepewność} -- wykrycie w strefie radarowej jest
        zdarzeniem losowym, którego prawdopodobieństwo nie jest znane agentowi.
  \item \textbf{Występuje konflikt celów} -- krótsza trasa przez strefę
        radarową oznacza mniejsze zużycie energii, ale większe ryzyko;
        objazd jest bezpieczniejszy, ale może wymagać dodatkowego ładowania.
  \item \textbf{Nagroda jest opóźniona} -- główna nagroda za sukces pojawia
        się dopiero po powrocie do bazy, więc agent musi przypisać zasługę
        wcześniejszym decyzjom (np. odwiedzeniu stacji ładowania).
\end{enumerate}

Te cechy odpowiadają klasycznemu sformułowaniu problemu decyzyjnego
Markowa~\cite{sutton2018}, w którym agent uczy się polityki maksymalizującej
oczekiwany zdyskontowany zwrot bez znajomości modelu przejść.

\section{Formalizacja RL/MDP}
\label{sec:formalizacja}

Problem modelujemy jako epizodyczny proces decyzyjny Markowa
$(\mathcal{S}, \mathcal{A}, P, R, \gamma)$.

\subsection{Agent i środowisko}

\begin{itemize}
  \item \textbf{Agent} -- dron; w każdym kroku wybiera jedną akcję.
  \item \textbf{Środowisko} -- mapa $W \times H$ wraz z bazą, celami,
        stacjami ładowania, przeszkodami, strefami radarowymi, patrolem
        i korytarzem zamykanym. Środowisko przechowuje bieżący stan, wykonuje
        przejścia i zwraca nagrody.
\end{itemize}

\subsection{Stan}

Stan agenta to krotka
\begin{equation}
  s = (x,\; y,\; e,\; m,\; \varphi),
\end{equation}
gdzie:
\begin{itemize}
  \item $(x, y)$ -- pozycja drona na mapie,
  \item $e \in \{0, 1, \dots, E_{\max}\}$ -- poziom baterii,
  \item $m \in \{0,1\}^K$ -- maska bitowa sfotografowanych celów,
  \item $\varphi \in \{0, \dots, T_p - 1\}$ -- faza patrolu
        ($\varphi = t \bmod T_p$, gdzie $T_p$ to okres trasy patrolu).
\end{itemize}

\paragraph{Dlaczego faza patrolu jest w stanie?} Pozycja patrolu zależy od
czasu. Gdyby $\varphi$ nie było częścią stanu, ten sam stan $(x,y,e,m)$
prowadziłby do różnych skutków w zależności od chwili, w której agent się
w nim znalazł -- własność Markowa nie byłaby spełniona. Pozycję patrolu da się
odtworzyć z $\varphi$, więc nie trzeba jej przechowywać osobno.

\paragraph{Dlaczego stan korytarza $Z$ nie jest w stanie?} Korytarz jest
zamknięty dokładnie wtedy, gdy sfotografowano cel~1, czyli gdy $m_1 = 1$.
Informacja ta jest już zawarta w masce $m$. Gdyby zdarzenie zależało od innego
warunku (np. czasu), stan należałoby rozszerzyć o dodatkową flagę.

\paragraph{Rozmiar przestrzeni stanów.} Górne oszacowanie wynosi
\begin{equation}
  |\mathcal{S}| \le W \cdot H \cdot (E_{\max}+1) \cdot 2^K \cdot T_p .
\end{equation}
Dla mapy \emph{easy} ($10\times10$, $E_{\max}=30$, $K=3$, $T_p=6$) daje to
$100 \cdot 31 \cdot 8 \cdot 6 = 148\,800$ stanów, a po uwzględnieniu
$|\mathcal{A}| = 5$ około $7{,}4 \cdot 10^5$ wartości $Q(s,a)$. Taka tablica
mieści się w pamięci (kilka MB), więc podejście tablicowe jest wykonalne.
Faktycznie osiągalnych stanów jest mniej, ponieważ pola z przeszkodami
nie są dostępne.

\subsection{Akcje}

\[
  \mathcal{A} = \{\textsc{góra},\ \textsc{dół},\ \textsc{lewo},\ \textsc{prawo},\ \textsc{czekaj}\}.
\]

\begin{itemize}
  \item Ruch na pole poza mapą, na przeszkodę lub na zamknięty korytarz jest
        \textbf{akcją niepoprawną}: dron pozostaje na miejscu, zużywa energię
        i otrzymuje dodatkową karę.
  \item \textsc{czekaj} na stacji ładowania zwiększa energię o $\Delta_L$
        (do $E_{\max}$). Poza stacją oznacza zawis i zużywa energię jak ruch.
        Akcja ta pozwala też przeczekać przejście patrolu.
\end{itemize}

\subsection{Przejście (pojedynczy krok interakcji)}

W kroku $t$ środowisko wykonuje kolejno:
\begin{enumerate}
  \item wyznaczenie nowej pozycji na podstawie akcji $a_t$ (z obsługą akcji
        niepoprawnych),
  \item aktualizację energii: $e \leftarrow e - 1$ albo
        $e \leftarrow \min(e + \Delta_L, E_{\max})$ przy ładowaniu,
  \item przesunięcie patrolu: $\varphi \leftarrow (\varphi + 1) \bmod T_p$,
  \item sprawdzenie kolizji z patrolem (również „minięcia się'' na sąsiednich
        polach, gdy dron i patrol zamieniają się miejscami),
  \item jeżeli dron jest na niesfotografowanym celu $i$ -- ustawienie $m_i = 1$
        (dla $i=1$ zamknięcie korytarza $Z$),
  \item jeżeli dron jest w strefie radarowej -- losowanie wykrycia
        z prawdopodobieństwem $p_{\text{det}}$,
  \item sprawdzenie warunków zakończenia i zwrócenie $(s_{t+1}, r_{t+1},
        \textit{terminated}, \textit{truncated})$.
\end{enumerate}

Interfejs środowiska będzie zgodny z biblioteką Gymnasium~\cite{gymnasium}
(\texttt{reset()}, \texttt{step(a)}).

\subsection{Funkcja nagrody}

\begin{table}[h]
\centering
\begin{tabular}{lrl}
  \toprule
  Zdarzenie & Nagroda & Uzasadnienie \\
  \midrule
  każdy krok & $-1$ & zachęca do krótkich tras \\
  akcja niepoprawna & $-2$ (dodatkowo) & zniechęca do „odbijania się'' od ścian \\
  sfotografowanie nowego celu & $+10$ & nagroda pośrednia, ułatwia eksplorację \\
  powrót do bazy ze wszystkimi celami & $+50$ & główny cel misji \\
  wykrycie, kolizja z patrolem & $-50$ & porażka misji \\
  rozładowanie baterii & $-50$ & porażka misji \\
  \bottomrule
\end{tabular}
\caption{Wstępna funkcja nagrody (wartości do strojenia po pierwszych testach).}
\label{tab:reward}
\end{table}

Kara za porażkę ($-50$) jest celowo równa nagrodzie za sukces, aby agent nie
wolał szybko zakończyć epizodu porażką zamiast długo latać z karą $-1$ za krok.
Przy długości epizodu rzędu 40--60 kroków suma kar za kroki jest porównywalna
z karą za porażkę -- ten balans trzeba sprawdzić eksperymentalnie
(zob. sekcja~\ref{sec:otwarte}).

\subsection{Warunki zakończenia epizodu}

\begin{itemize}
  \item \textbf{Sukces} (\textit{terminated}): dron jest w bazie i $m = (1,\dots,1)$.
  \item \textbf{Porażka} (\textit{terminated}): wykrycie w strefie radarowej,
        kolizja z patrolem lub $e = 0$ poza bazą.
  \item \textbf{Przekroczenie limitu} (\textit{truncated}): $t \ge T_{\max}$.
\end{itemize}

Rozróżnienie \textit{terminated}/\textit{truncated} jest istotne przy
aktualizacji wartości: po zakończeniu właściwym wartość stanu następnego
wynosi $0$, natomiast po obcięciu limitem kroków agent nadal powinien
korzystać z oszacowania $\max_a Q(s_{t+1},a)$.

\subsection{Ograniczenia środowiska}

\begin{itemize}
  \item energia $E_{\max}$ (wstępnie 30) i przyrost ładowania $\Delta_L$
        (wstępnie 5 na krok),
  \item limit kroków $T_{\max}$ (wstępnie 200),
  \item prawdopodobieństwo wykrycia $p_{\text{det}}$ (wstępnie 0{,}3 na krok),
  \item brak możliwości wejścia na przeszkody i zamknięty korytarz.
\end{itemize}

\subsection{Współczynnik dyskontowania}

Wstępnie $\gamma = 0{,}99$. Główna nagroda pojawia się dopiero na końcu
długiego epizodu, więc zbyt mała wartość $\gamma$ sprawiłaby, że agent
praktycznie nie odczuwałby nagrody za powrót do bazy. Wpływ $\gamma$ zostanie
sprawdzony w eksperymentach.

\subsection{Miary sukcesu projektu}

\begin{itemize}
  \item odsetek misji zakończonych sukcesem (w fazie ewaluacji, bez eksploracji),
  \item średni zwrot z epizodu,
  \item średnia długość udanego epizodu,
  \item odsetek porażek z podziałem na przyczyny (wykrycie, patrol, bateria),
  \item stabilność wyników między uruchomieniami z różnym ziarnem losowym.
\end{itemize}

Projekt uznajemy za udany, jeżeli agent uczący się istotnie przewyższa agentów
bazowych na wszystkich mapach, a zachowanie agenta (np. omijanie radarów,
planowanie ładowania, kolejność celów) daje się zinterpretować.

\subsection{Elementy wymagające weryfikacji po pierwszych testach}
\label{sec:otwarte}

\begin{itemize}
  \item wartości nagród -- czy agent nie uczy się „samobójczych'' polityk
        albo krążenia w miejscu,
  \item czy nagroda $+10$ za cel nie powoduje, że agent zbiera cele i nie
        wraca do bazy,
  \item wartość $p_{\text{det}}$ -- zbyt duża sprawi, że radary staną się
        zwykłymi ścianami, zbyt mała -- że będą bez znaczenia,
  \item czy osiągalna przestrzeń stanów jest wystarczająco mała dla
        tablicowego Q-learningu na mapie \emph{hard},
  \item ewentualne dodanie losowego znoszenia drona przez wiatr
        (dodatkowe źródło niepewności).
\end{itemize}

% Etap E1 (propozycja tematu): oddawane są rozdziały 1-3 oraz sekcja Wsparcie AI.
% Kolejne rozdziały odkomentować w etapie E2.
% \section{Opis środowiska}
\label{sec:srodowisko}

\todo{Etap E2: opis implementacji środowiska (klasa, format plików map,
\texttt{reset}/\texttt{step}, obsługa ziarna losowego), opis map
\emph{easy}/\emph{medium}/\emph{hard} oraz przetestowanych przypadków brzegowych
(wyjście poza mapę, brak energii, wejście w radar, kolizja z patrolem,
zamknięcie korytarza).}

% \section{Agent bazowy}
\label{sec:baseline}

Planowane są dwa agenty bazowe:
\begin{itemize}
  \item \textbf{agent losowy} -- wybiera akcję z rozkładu jednostajnego;
        wyznacza dolną granicę wyników,
  \item \textbf{agent zachłanny} -- wyznacza najkrótszą ścieżkę (BFS) do
        najbliższego niesfotografowanego celu, a po zebraniu wszystkich celów
        -- do bazy; ignoruje baterię, radary i patrol. Jest to uczciwy
        punkt odniesienia, ponieważ reprezentuje „rozsądną'' strategię
        bez uczenia.
\end{itemize}

\todo{Etap E2: wyniki agentów bazowych na każdej mapie.}

% \section{Zastosowany algorytm}
\label{sec:algorytm}

Planowane jest porównanie metod tablicowych~\cite{sutton2018}:
Q-learning (off-policy), SARSA (on-policy) oraz Q($\lambda$).
Spodziewana różnica: SARSA, uwzględniając eksplorację w swojej polityce,
powinna wybierać trasy dalej od stref radarowych niż Q-learning
(analogicznie do przykładu \emph{cliff walking}).

\todo{Etap E3: reprezentacja tablicy $Q$ (indeksowanie stanu), strategia
$\varepsilon$-zachłanna i jej harmonogram, hiperparametry, rozdzielenie
fazy uczenia i ewaluacji.}

% \section{Plan eksperymentów}
\label{sec:plan}

\todo{Etap E2: porównywane warianty (algorytmy, mapy, wartości $\alpha$,
$\gamma$, $\varepsilon$, warianty funkcji nagrody), liczba epizodów,
liczba uruchomień z różnymi ziarnami, lista planowanych wykresów i tabel.}

% \section{Wyniki i analiza}
\label{sec:wyniki}

\todo{Etap E3: krzywe uczenia, porównanie z agentami bazowymi, analiza
nieudanych epizodów, wpływ hiperparametrów, wnioski.}

% \section{Ograniczenia rozwiązania}
\label{sec:ograniczenia}

\todo{Etap E3: np. pełna obserwowalność stanu, skalowanie tablicy $Q$
z rozmiarem mapy, deterministyczny patrol, wrażliwość na dobór nagród.}

% \section{Instrukcja uruchomienia}
\label{sec:uruchomienie}

\todo{Uzupełniać na bieżąco, zgodnie z plikiem \texttt{README.md}.}

\section{Wsparcie AI}
\label{sec:ai}


\bibliographystyle{plain}
\bibliography{references}

\end{document}
